\chapter{如何寫一本書：一位暢銷作家的 13 個步驟}

這場講座以一連串刻意安排的限制與釋放展開。我們從講者開始的地方開始：先反對網路上那種輕易的承諾，彷彿一本書可以靠某種小技巧迅速製造出來。然後，一步一步地，我們從工作條件走到規模，從規模走到結構，從結構走到讀者的前進感，最後再從起草時的自由，走向毫不留情的修訂。最終浮現出來的，不是一袋零散建議，而是一套關於長篇手稿如何被建造起來的有紀律模型。

\section{品質優先：空間、工具，以及拒絕找藉口}

講座以一種論戰式姿態開場。這個世界上有太多關於速度的承諾、太多關於立刻成功的公式、太多關於如何快速造出一本暢銷書的濃縮配方。講者立刻拒絕了這種氣氛。他不承諾暢銷，也不承諾速度；他承諾的是十三個基礎步驟。這裡的主導性對比，可以簡潔地寫成
\begin{equation}
\text{速度} \neq \text{目標}, \qquad \text{品質} = \text{目標}.
\end{equation}
這條關係決定了後面一切的語氣。既然品質才是目標，那麼我們首先關心的就不是加速，而是什麼樣的條件能使長期工作真正成為可能。

也正因此，講座才從如此具體的地方開始。我們先建立一個寫作空間，接著再把寫作工具準備齊全。這些最初的步驟並不華麗，但它們的目的，是清除初學者最愛使用的那種藉口。講者說得很明白：我們不應該說自己沒有地方寫。幾把椅子架上一塊木板，曾經就能當桌子；餐廳可以用；車子可以用；吵雜的新聞編輯室也曾經可以用。重點不是說理想條件無足輕重，而是說：工作可以在那些僅僅「可用」的條件下開始。

同樣的嚴肅也適用於工具。我們不需要一整套花俏的設備，但我們確實需要足夠穩定的條件，讓自己不必每隔幾分鐘就被迫中斷專注。一台不與我們作對的電腦、一張不折磨背與脖子的椅子，以及那些不然就會讓我們滿屋子找來找去的基本物件，全部都服從於同一個邏輯：在任何大工程真正開始移動之前，我們先把摩擦降到最低。

\subsection{Question \& Answer}

\paragraph{Question.} 我們是不是一定要等到理想條件出現，才能開始？

\paragraph{Answer.} 不。講座區分的是理想條件與可用條件。隱私、安靜與更好的設備確實有幫助，但它們的缺席並不能成為拖延的理由。我們真正需要的，是一個能夠反覆出現的空間與配置。關鍵詞不是完美，而是已經建立起來。

\section{拆解問題的規模，並測試這個點子是否夠大}

當局部條件確定之後，講座立刻抬高下一個障礙。一本書，單單憑藉它的體積，就足以讓人退縮。講者對此的回答不是靈感，而是換一個尺度來思考。一部書不能作為一個五百頁的物件直接進攻，而應被拆成小單位：
\begin{equation}
\text{一本書} \to \text{章節} \to \text{段落} \to \text{句子}.
\end{equation}
這正是那個「吃大象」類比的實際內容。我們不是一口吞下整體，而是吃掉眼前這一口。

同樣的點，也可以再更示意性地寫成：
\begin{equation}
\text{龐大任務} \Rightarrow \text{小型可行動單位}.
\end{equation}
講座在這裡停留得特別久，因為這是第一個重要的心理反轉。書依然很大，但作者下一步要做的事情卻已經變得很小。

然而，到了這裡，講者又拒絕讓我們過早得到一個過於容易的結論。並不是每一個計畫都值得被毅力拯救。有些點子本身太小，承受不起它被要求佔據的形式。因此，第四步就來了：先確定你有一個夠大的點子。這個區分並不含糊。小概念，屬於部落格或文章；大概念，才屬於一本到書。
\begin{equation}
\text{小概念} \Rightarrow \text{部落格／文章}, \qquad
\text{大概念} \Rightarrow \text{一本書}.
\end{equation}
講者不是抽象地談這件事，而是直接提名字來讓它變得可見。談非虛構時，他舉 \textit{How to Win Friends and Influence People}；談虛構時，他舉 \textit{Harry Potter}。重點不是模仿，而是規模。

\subsection{Question \& Answer}

\paragraph{Question.} 我們怎麼知道一個點子是否真的長得出一本到書？

\paragraph{Answer.} 講座給出的是一種非常實用的測試：把這個點子講給另一個人聽。若在複述過程中，它一直變大、一直打開，那它多半具備足夠的規模；若它一再撞上同一面牆，那問題就不是你的耐力，而是點子本身的尺寸不夠。

\section{提綱、虛構與非虛構，以及中段馬拉松}

到了第五步，講座進入了它的第一個結構高峰。白板之所以重要，正因為它記錄了這場演講在某一時刻如何從一堆實用規則，凝結成一套形狀理論。

\begin{figure}[h]
\centering
\includegraphics[width=0.88\textwidth]{lecture_05_figure_03.png}
\caption{第五步的白板提綱：虛構、非虛構、中段，以及提綱的角色。}
\end{figure}

若把白板整理成乾淨的排版，便是
\begin{align}
5.\ \text{提綱} \notag\\
\text{虛構 vs.\ 非虛構}\qquad (\text{King})\qquad (\text{經紀人／編輯}) \notag\\
\text{中段馬拉松} \notag\\
\text{同一種結構} \notag\\
\text{提綱是為了 } \underline{\text{你}} \text{ 服務。} \notag
\end{align}

講者先從觀眾必定會感受到的一種張力開始。既然有些虛構作家是 pantser，靠邊寫邊發現而不是靠詳細規劃，那為什麼還要堅持提綱？他的回答很仔細。提綱的程度當然可以不同，但方向感不能消失。Stephen King 代表其中一端：把有趣的角色放進困難情境裡，然後透過書寫去發現會發生什麼。即便如此，某種方向感仍然是必要的。非虛構則位於另一端。在那裡，作者要對經紀人與編輯負責，他們會想知道你要說什麼、怎麼說、資料來源在哪裡、論點將會往哪裡走。

也正因如此，白板上的下一句話才會如此重要：\textit{same structure}。不同文類的表面資源當然有別，但底層運動其實相同。講座把這條共同法則說得非常簡潔：
\begin{equation}
\text{鋪陳} \to \text{回收}.
\end{equation}
這就是本章真正的數學脊骨。虛構經常透過場景、對話、反轉與衝突來實現它；非虛構則可能透過問題、方法、說明與結果來實現它。但在兩種情況下，讀者都一樣是被「先被種下的期待」與「後來被完成的滿足」所帶動。

\begin{definition}
所謂中段馬拉松，是指手稿大約位於中點與四分之三處之間的那段區域，在那裡，作者最容易感受到篇幅的拖力，並開始想要灌水、棄筆，或只是硬撐。
\end{definition}

若用一種緊湊的編者速記，我們可以寫成
\begin{equation}
x \in \left[\tfrac12,\tfrac34\right] \Rightarrow \text{結構疲勞的風險提高},
\end{equation}
同時記得，這裡的 \(x\) 是我們自己的摘要，而不是講座裡使用的符號。

講者在這裡最強的一條主張是：我們不能只是熬過中段，而必須在中段裡真正活起來。若中段對我們而言已經顯得乏味，那讀者其實早就睡著了。這條原則對非虛構同樣成立。一本關於製作模型船的非虛構書，不應該只是列出步驟而已；它也應該設計困難，讓最後的解決顯得是掙來的。

\subsection{Question \& Answer}

\paragraph{Question.} 如果我們是 pantser，為什麼還需要提綱？

\paragraph{Answer.} 因為發現不等於漫無方向。虛構作者可以用比較輕的提綱，非虛構作者則可以用比較重的提綱，但在兩種情況下，提綱都為長篇形式提供了一條方向，使它不至於在前幾頁過後鬆散垮掉。講座最後的修正同樣重要：提綱是一件工具，它是為了服務書；而書又是為了服務作者。如果發展中的手稿被證明比原本的方案更好，那就應該修改提綱，而不是反過來讓書去屈從提綱。

\section{進度表、截止日，以及作為控制變數的研究}

當結構被確立之後，講座重新回到執行。第六步，是把寫作意志轉換成算術。作者估計整本書的總量，再估計可用天數，然後相除：
\begin{equation}
\text{每日配額} = \frac{\text{目標頁數或字數}}{\text{可用天數}}.
\end{equation}
這條公式的重要性，不在數學上的優美，而在它帶來的道德清晰感。只有當截止日被換算成每天的義務時，它才真正變成一件真實的事。

講座在調整這件事上也同樣實際。我們也許一開始估出每天十頁，然後很快發現，四頁或五頁才更貼近現實。這時，我們當然可以修改進度表。但一旦修改完並定下來，它就必須被當成神聖之物。最容易犯的錯，是以為寫作時間會偶然出現。講座講得正好相反：你必須把一天裡別的東西犧牲掉。

\paragraph{Worked example.} 講者用自己的進度舉了一個很具體的例子：
\begin{align}
70{,}000 - 67{,}500 &= 2{,}500,\\
\text{明天的配額} &= 2{,}500\ \text{字}.
\end{align}
這個算術本身很簡單，但它的意義並不簡單。它是把一種模糊的野心，轉化成一條銳利要求的動作。

接著，講者又補上一條讓人清醒的出版經驗法則：
\begin{equation}
\text{能守住截止日的人} \approx \frac{1}{100}.
\end{equation}
這當然不是定理，而是一種專業估計，用來提高我們對嚴肅性的標準。光是能夠守住一個截止日，本身就已經把你和大多數寫作者區分開了。

第七步，便把討論從進度表擴展到知識。研究在非虛構中當然是必要的，但講者也堅持：許多作者都低估了它在虛構寫作中的重要性。歷史、飛機、武器，只要錯一個細節，真正懂行的讀者就會立刻察覺。講座在這裡採用的隱喻非常準確：
\begin{equation}
\text{研究} \neq \text{主菜}, \qquad
\text{研究} = \text{調味}.
\end{equation}
這個隱喻很好地維持了層級。研究提供的是特異性與可信度，而不是對故事本身的取代。講者隨即又舉出自己工作中的一個具體例子：他正在使用一張從公元前 2000 年延伸到當代的世界歷史圖表，以精確追蹤人物的出生、死亡與交疊關係。研究的存在，是為了支撐敘事線，而不是把敘事淹沒掉。

\subsection{Question \& Answer}

\paragraph{Question.} 要做到什麼程度的研究，才算足夠，而又不至於讓研究本身反過來佔據整本書？

\paragraph{Answer.} 做到足以使作品變得具體、可信、技術上站得住腳，卻又不至於讓讀者被迫欣賞作者的資料卡片。當研究能夠把頁面磨利，卻自己不變成頁面主題時，它就完成了自己的工作。

\section{讀者優先的開場，以及張力的邏輯}

當結構、進度表與研究都就位之後，講座便果斷地轉向讀者。第八步從開場句開始。這個重點不在於「第一句是魔法」，而在於：第一句會確定語調，並建立壓力。講者把開場句分成幾個類型：
\begin{equation}
\{\text{令人驚奇的},\ \text{戲劇性的},\ \text{哲理性的},\ \text{詩性的}\}.
\end{equation}
接著，他舉出例子。Orwell 在 \textit{1984} 的開場，被用來說明一種立刻讓世界失去穩定的驚奇。Elizabeth Gilbert 的非虛構例子，則是靠銳利而具體的動作做到同樣的事。Toni Morrison 在 \textit{Paradise} 裡那句著名開頭，則被歸到戲劇性陳述一類。講座並沒有長久停留在例子的層次，而是立刻把開頭句的局部問題，轉化成一條整體規則。

這條規則就是「讀者優先」的篩選器：
\begin{equation}
\text{一個決定} \xrightarrow{\text{讀者優先}} \{\text{保留},\ \text{修訂},\ \text{刪除}\}.
\end{equation}
講者直接這麼說：手稿中的每一個決定，都應當先經過一個問題——它對讀者有什麼作用？不是作者優先，不是編輯優先，也不是評論家優先，而是讀者優先。這是整場講座中最強的一次概念轉向之一。開場之所以重要，不是因為它可以炫技，而是因為它教會我們之後所有決策都該如何被裁判。

第九步則把這個原則，從開場句擴展到整部手稿的生命。究竟是什麼讓讀者持續翻頁？講座的回答是：壓力。
\begin{equation}
\text{衝突／張力} \Rightarrow \text{翻頁的動量}.
\end{equation}
對虛構而言，這個說明非常鮮明。和氣融融在生活裡也許很好，但放到頁面上往往就死了。場景真正活起來，是當某種埋藏著的緊繃浮現出來、當某個角色說出意料之外的話、當表面的對話已經無法自己解釋自己時。

但非虛構並不因此豁免。講者明白問道：那麼在非虛構裡，張力是怎麼來的？他的回答不是戲劇性的，而是結構性的：
\begin{equation}
\text{承諾} \to \text{兌現}.
\end{equation}
一本好的非虛構書，會在開頭幾頁就向讀者承諾：到最後，你將會得到某種東西；而接下來它就必須真的把那件東西交出來。在這個語境裡，張力不是對抗，而是被延後的滿足。

\subsection{Question \& Answer}

\paragraph{Question.} 沒有反派、戰鬥或通俗戲劇的情況下，非虛構要如何創造張力？

\paragraph{Answer.} 方法是給讀者一件值得等待的事。我們先提出一個困難、一個不可能、或一種想要實現的轉變，接著再把整部手稿組織成：後面的章節必須真正完成前面的章節所承諾的東西。在非虛構裡，張力常常就是「宣布」與「抵達」之間被嚴格控制的距離。

\section{起草的自由、中段的再臨，以及擲地有聲的結尾}

第十步改變了時間尺度。到目前為止，講座大多數時候都是從外部描述結構；到了這裡，它開始問：在真正寫第一稿的那個時刻，作者應當允許什麼發生？答案同時嚴厲又解放：把內在編輯器關掉。
\begin{equation}
\text{關掉編輯器來起草} \to \text{之後再回來編輯} \to \text{用刪減來增加力量}.
\end{equation}
這個次序非常重要。第一稿不是每一個詞都要即時審判的地方。講者用很個人的方式說明這件事：內在批評者會坐在肩膀上，對每一件事提出異議。而在起草階段，這個批評者必須被靜音。第一稿的全部目的，只是先把材料放下來。

\subsection{Question \& Answer}

\paragraph{Question.} 內在編輯器什麼時候應該閉嘴？又什麼時候應該回來？

\paragraph{Answer.} 它在起草的時候應該閉嘴，因為過早的批評會摧毀整體運動；而只有在過了一段時間之後，當作者可以把手稿看成某個有待塑形的東西，而不再是某個尚在掙扎著誕生的東西時，它才應該回來。

就在這裡，講座又做了一次它最重要的回顧。第十一步明白地把話題帶回「中段馬拉松」。這不是為了重複而重複，而是在那個早先被提出的結構原則，此刻作為一種反覆發生的真實危機再次現身。講者坦白說，這正是他自己最想放棄的地方。真正的危險是灌水：在虛構裡加入多餘場景，在非虛構裡加入多餘論點，讓篇幅延長，卻不增加壓力。

對策依舊是：不要只是忍過去，而要在其中活得更強。中段仍然必須維持真實的鋪陳與回收。在講者自己正在進行的作品裡，一種讓中段不死的方法，就是在公元前 2000 年、當代，以及越戰之間來回切換，如此一來，手稿能不斷更新自己的壓力，而不至於被壓平成一整帶解說。

接著，第十二步轉向收束。結尾不能無聲無息地洩掉，它必須落下，像劇院裡幕布放下來一樣有力。若存在幾種可能的結尾，講座的規則非常清楚：
\begin{equation}
\text{候選結尾} \to \text{選出最有情感力量的那一個}.
\end{equation}
這並不等於鼓勵煽情，而是意味著：讀者記住的是那些真正打動他們的東西，因此書的最後一個動作，必須回收那種一路被準備出來的力量。

\section{凶狠的自我編輯與專業判斷}

講座的最後一個段落，又把標準再收緊了一次。我們被要求成為凶狠的自我編輯者。這裡的「凶狠」不是神祕語，而是字面上的：要有攻擊性地剪、精修，讓手稿開始歌唱。如果我們不自己做這些工作，那麼前面的一切就已經部分浪費了。

接著，講者引入了手稿實際上會如何被閱讀的那種專業壓力。經紀人與編輯要讀的東西太多，時間卻太少。很多時候，他們在兩分鐘內、或者在一兩頁之內，就已經能判斷這份手稿值不值得繼續讀下去。這裡的邏輯是殘酷而簡單的。競爭極其龐大，因此開頭幾頁必須已經被磨到銳利。自我編輯因此絕不是最後多餘的一道裝飾工序，而是手稿被真正準備好、得以不帶歉意地去面對世界的那個步驟。

在結尾處，講座又回到了自己的起點。它一開始反對捷徑與速度炒作，最後則展示：真正站在那些東西對面的，是工作條件、一個夠大的點子、一份提綱、一張進度表、研究、讀者優先的判斷、張力、一個自由的初稿、一個真正活起來的中段、一個能夠落地的結尾，以及一位足夠殘酷、肯一路刪到手稿真正帶出它全部力量的編輯者——而這位編輯，正是作者自己。

\section{Summary}

這場講座一步一步地向前推進，但它的力量其實來自於步驟之間如何彼此累積。我們先去除藉口，建立能夠工作的條件；再拆小問題的規模，然後檢驗這個概念是否大到值得被寫成一本到書。我們列提綱，不是因為要做提綱的奴隸，而是因為長篇形式需要方向。我們把野心換算成每日配額，把研究當成調味而不是炫耀，並把每一個局部決定都交給讀者的經驗來裁判。我們先起草，再編輯；我們不只是熬過中段，而是讓它活起來；我們選那個最能深深打動人的結尾；最後，再以凶狠的方式把整部書修到銳利。這場講座之所以嚴肅，建立在一個非常簡單的真理上：一本到書不是被願望喚出來的，而是被建造、被維持，並被磨利出來的。